\begin{example}
	Напишем явную формулу для вычисления обратного элемента к числу $a := a_0 + a_1 \theta$ в поле~$\QQ[\theta]$, где $\theta^2 + \theta + 1 = 0$.

	Это равносильно нахождению обратного элемента в поле
	\[
		\frac{\QQ[x]}{(x^2 + x + 1)}
		.
	\]

	Пусть $b := b_0 + b_1\theta$ --- обратный к~$a$. Тогда $ab = 1$, то есть
	\begin{align*}
		(a_0 + a_1 \theta) (b_0 + b_1 \theta) &= 1
		\\
		a_0 b_0 + (a_1 b_0 + a_0 b_1) \theta + a_1 b_1 \theta^2 &= 1
		.
	\end{align*}
	Так как $\theta^2 + \theta + 1 = 0$, то $\theta^2 = -\theta - 1$ (стершую степень многочлена, аннулирующего~$\theta$, можно выразить через его меньшие степени). Тогда
	\begin{align*}
		a_0 b_0 + (a_1 b_0 + a_0 b_1) \theta + a_1 b_1 (-\theta - 1) &= 1
		\\
		(a_0 b_0 - a_1 b_1) +
		(a_1 b_0 + a_0 b_1 - a_1 b_1) \theta &= 1
	\end{align*}
	Приравняем коэффициенты при одинаковых степенях~$\theta$:
	\[
		\begin{system}
			a_0 b_0 - a_1 b_1 &= 1,
			\\
			a_1 b_0 + (a_0 - a_1) b_1 &= 0.
		\end{system}
	\]
	В матричной записи
	\[
		\begin{pmatrix}
			a_0 & -a_1
			\\
			a_1 & a_0 - a_1
		\end{pmatrix}
		\Vect{b_0 \\ b_1}
		=
		\Vect{1 \\ 0}.
	\]
	Так как мы находимся в поле, то всякий ненулевой элемент обратим. Решим систему методом Крамера:
	\[
		\Delta :=
		\begin{determ}
			a_0 & -a_1
			\\
			a_1 & a_0 - a_1
		\end{determ}
		=
		a_0 (a_0 - a_1) + a_1^2 =
		a_0^2 - a_0 a_1 + a_1^2 =
		(a_0 - a_1)^2 + a_0 a_1.
	\]

	\textstars

	Есть другой способ.

	Достаточно найти обратные к элементам вида~$\theta - a$. Это можно сделать по алгоритму Евклида: если $h(\theta)$ --- класс, обратный к $\theta - a$, то $h(\theta)$ задаётся таким многочленом $h \in \QQ[x]$, что $[h(x)][x-a] = [1]$, то есть $[h(x)(x-a)] = [1]$, то есть $h(x)(x-a) + \bigl( (x^2 + x + 1) \bigr) = 1 + \bigl( (x^2 + x + 1) \bigr)$, то есть
	\[
		h(x) (x - a) + g(x)(x^2 + x + 1) = 1
	\]
	для некоторого $g \in \QQ[x]$.

	Такое разложение можно найти по алгоритму Евклида. Пусть $E_0 := x^2 + x + 1$, $E_1 := x-a$. Найдём $E_2 := E_0 \bmod E_1$ (поделим столбиком):
	\[
		x^2 + x + 1 =
		(x-a) \cdot x +
		\underbrace{%
			\bigl( (a+1)x + 1 \bigr)
		}_{E_2}
	\]
	Итак, $E_0 = xE_1 + E_2$, откуда $E_2 = E_0 - xE_1$; далее найдём $E_3 := E_1 \bmod E_2$:
	\[
		x-a = 
		\bigl( (a+1)x + 1 \bigr)
		\cdot
		\frac{1}{a+1}
		+
		\underbrace{%
			\bigl(
				-a - \frac{1}{a+1}
			\bigr)
		}_{E_3}
	\]
	Итак,
	\begin{align*}
		E_3 &=
		E_1 - \frac{E_2}{a + 1}
		\nlinea{=}
		E_1 - \frac{E_0 - xE_1}{a + 1}
		\nlinea{=}
		\frac{-E_0 + (a+1+x)E_1}{a + 1}
	\end{align*}
	Или, что то же самое,
	\begin{align*}
		-a - \frac{1}{a+1} &=
		\frac{-E_0 + (a+1+x)E_1}{a + 1}
		\\
		-a(a+1) - 1 &=
		-E_0 + (a+1+x)E_1
		\\
		1 &=
		\frac{-E_0}{-a^2 - a - 1} +
		\frac{(a+1+x)E_1}{-a^2 - a - 1}
		\\
		1 &=
		\frac{-E_0}{-a^2 - a - 1} +
		\frac{(a+1+x)E_1}{-a^2 - a - 1}
	\end{align*}

\end{example}

